\documentclass[12pt,a4paper]{article}

\usepackage[utf8]{inputenc}
\usepackage[T2A]{fontenc}
\usepackage[russian]{babel}
\usepackage{amsmath}
\usepackage{amssymb}
\usepackage{graphicx}
\usepackage{hyperref}
\usepackage[margin=2.5cm]{geometry}
\usepackage{amsopn}
\usepackage{textcomp}

\title{История развития небесной механики: от Птолемея до спутниковых группировок}
\author{}
\date{}

\begin{document}

\maketitle
\tableofcontents
\newpage

%=============================================================================
\section{Введение}
%=============================================================================

Небесная механика~--- одна из древнейших ветвей науки, зародившаяся из
стремления человека понять движение светил на ночном небе. Уже в Древней
Греции философы пытались построить стройную модель Вселенной, опираясь на
наблюдения и геометрию. Птолемей систематизировал эти идеи в геоцентрическую
систему мира, которая господствовала более тысячи лет.

Переворот произошёл в XVI--XVII веках, когда Николай Коперник предложил
гелиоцентрическую модель, Джордано Бруно развивал идею о бесконечности
Вселенной и множественности миров, а Галилео Галилей подтвердил
гелиоцентризм наблюдениями в телескоп. Одновременно Иоганн Кеплер, опираясь
на наблюдательные таблицы Тихо Браге, открыл три закона движения планет.

Исаак Ньютон совершил великий синтез: он объединил земную механику Галилея и
небесную механику Кеплера в единую теорию, основанную на законе всемирного
тяготения. Для этого ему пришлось создать дифференциальное и интегральное
исчисление. Вклад Эйлера, Лагранжа и Лапласа в XVIII--XIX веках превратил
небесную механику в одну из самых развитых областей математической физики.

В XX веке Альберт Эйнштейн показал, что ньютоновская теория тяготения~---
лишь приближение к более глубокой общей теории относительности, а Анри
Пуанкаре открыл хаотическое поведение в задаче трёх тел, заложив основы
теории динамических систем.

Сегодня небесная механика переживает новый расцвет: если раньше учёные
изучали движение естественных небесных тел, то теперь перед нами стоит задача
создания и управления искусственными спутниками, что порождает целый новый
пласт математики~--- теорию управления.

%=============================================================================
\section{Система мира Птолемея}
%=============================================================================

Клавдий Птолемей (ок.~100--170~н.\,э.) создал грандиозный труд <<Альмагест>>,
в котором описал геоцентрическую модель Вселенной. В центре мира находилась
неподвижная Земля, а все небесные тела~--- Солнце, Луна, планеты и звёзды~---
обращались вокруг неё.

Для объяснения наблюдаемого движения планет, включая попятные петли,
Птолемей использовал систему \textbf{деферентов} и \textbf{эпициклов}.
Деферент~--- большая окружность с центром вблизи Земли, по которой движется
центр малой окружности~--- эпицикла. Планета движется по эпициклу, и
суперпозиция двух круговых движений порождает сложную кривую, которая
качественно воспроизводит наблюдаемую траекторию планеты на небосводе.

Замечательный математический факт состоит в том, что, добавляя всё больше
эпициклов (эпициклы на эпициклах), можно приблизить \emph{любую} замкнутую
кривую с произвольной точностью. Это непосредственно связано с разложением
в ряд Фурье. Действительно, траекторию на комплексной плоскости можно
записать как
\begin{equation}
  z(t) = \sum_{n=-N}^{N} c_n \, e^{i n \omega t},
\end{equation}
где каждое слагаемое $c_n e^{in\omega t}$ соответствует отдельному эпициклу
с радиусом $|c_n|$ и угловой скоростью $n\omega$. Таким образом, идея
Птолемея, хотя и основанная на ошибочной физической картине, фактически
предвосхитила одно из важнейших открытий математики~--- разложение функций
в ряд Фурье.

Для иллюстрации этой связи мы подготовили демонстрацию: профиль Птолемея
рисуется как кривая на комплексной плоскости с помощью набора вращающихся
эпициклов~--- каждый из них соответствует отдельной гармонике ряда Фурье.
Чем больше эпициклов мы берём, тем точнее воспроизводится контур.

Система Птолемея была практически полезной: она позволяла предсказывать
положения планет с приемлемой для того времени точностью и использовалась
астрономами и навигаторами на протяжении почти полутора тысячелетий. Однако
с накоплением более точных наблюдательных данных система требовала всё
больше эпициклов и эквантов, становясь громоздкой и неэлегантной. Это
подготовило почву для революционных идей Коперника.

%=============================================================================
\section{Гелиоцентрическая система Коперника}
%=============================================================================

В 1543 году Николай Коперник опубликовал труд <<De Revolutionibus Orbium
Coelestium>> (<<О вращениях небесных сфер>>), в котором поместил Солнце в
центр мира, а Землю превратил в одну из планет, обращающихся вокруг него.
Это была революционная идея, ломавшая тысячелетнюю традицию.

Коперник показал, что попятные движения планет, для объяснения которых
Птолемей строил сложную систему эпициклов, легко объясняются тем, что мы
наблюдаем за внешними планетами с движущейся Земли. Когда Земля обгоняет
внешнюю планету, та кажется движущейся назад на фоне звёзд.

Однако гелиоцентрическая идея столкнулась с мощным сопротивлением. Джордано
Бруно, пошедший ещё дальше и утверждавший, что Вселенная бесконечна, а
звёзды~--- такие же солнца с собственными планетными системами, был обвинён
в ереси и сожжён на костре в 1600 году. Его судьба стала символом конфликта
между наукой и догмой.

Решающие наблюдательные подтверждения гелиоцентризма пришли от Галилео
Галилея. В 1609--1610 годах, направив усовершенствованный телескоп на небо,
Галилей совершил ряд открытий:
\begin{itemize}
  \item \textbf{Фазы Венеры.} Галилей обнаружил, что Венера демонстрирует
    полный цикл фаз, подобно Луне. Это было возможно только в том случае,
    если Венера обращается вокруг Солнца, а не вокруг Земли.
  \item \textbf{Спутники Юпитера.} Галилей открыл четыре крупных спутника
    Юпитера (Ио, Европу, Ганимед и Каллисто), показав, что не всё во
    Вселенной вращается вокруг Земли. Существование другого центра
    обращения опровергало фундаментальный постулат геоцентризма.
  \item \textbf{Горы на Луне и пятна на Солнце.} Эти наблюдения разрушили
    представление о <<совершенстве>> небесных тел, которое было одним из
    столпов аристотелевской космологии.
\end{itemize}

Так был совершён один из величайших мировоззренческих переворотов в истории
науки~--- переход от геоцентрической к гелиоцентрической картине мира.

%=============================================================================
\section{Галилей и параболическое движение}
%=============================================================================

Помимо астрономических открытий, Галилео Галилей внёс фундаментальный вклад
в механику, заложив основы кинематики. Он первым систематически изучил
движение тел вблизи поверхности Земли и установил, что свободно брошенное
тело движется по параболе.

Главная экспериментальная трудность состояла в измерении коротких
промежутков времени~--- точных часов в XVII веке не существовало. Галилей
нашёл остроумное решение: он использовал собственный пульс как
приблизительный хронометр, а также проводил эксперименты с наклонной
плоскостью, на которой движение замедлялось и его было легче наблюдать.

Скатывая шарики по наклонной плоскости и измеряя пройденные расстояния
за равные промежутки времени, Галилей установил, что расстояние
пропорционально квадрату времени~--- тем самым открыв равноускоренное
движение.

Для тела, брошенного под углом $\theta$ к горизонту с начальной скоростью
$v_0$, уравнения движения записываются в виде:
\begin{align}
  x(t) &= v_0 \cos\theta \cdot t, \\
  y(t) &= v_0 \sin\theta \cdot t - \frac{1}{2} g t^2,
\end{align}
где $g$~--- ускорение свободного падения. Исключая время, получаем уравнение
параболы:
\begin{equation}
  y = x \tan\theta - \frac{g \, x^2}{2 v_0^2 \cos^2\theta}.
\end{equation}

Дальность полёта определяется формулой:
\begin{equation}
  L = \frac{v_0^2 \sin 2\theta}{g},
\end{equation}
и максимальна при $\theta = 45\textdegree$. Этот результат имел непосредственное
практическое значение для артиллерии.

Работы Галилея по земной механике стали одним из двух фундаментов, на
которых Ньютон впоследствии воздвиг единую механику. Вторым фундаментом
были небесные законы Кеплера.

%=============================================================================
\section{Законы Кеплера}
%=============================================================================

Иоганн Кеплер (1571--1630) совершил одно из величайших открытий в истории
науки~--- установил три закона движения планет. Основой для его работы
послужили многолетние, исключительно точные наблюдения Тихо Браге, который
на протяжении двадцати лет с беспрецедентной аккуратностью фиксировал
положения планет.

Кеплер получил доступ к таблицам Тихо Браге после его смерти в 1601 году
и потратил годы упорных вычислений, пытаясь согласовать наблюдения Марса
с различными моделями. Именно Марс, обладающий заметным эксцентриситетом
орбиты, позволил Кеплеру обнаружить отклонение от круговых орбит.

\subsection*{Первый закон}
Каждая планета движется по \textbf{эллипсу}, в одном из фокусов которого
находится Солнце. Уравнение эллипса в полярных координатах:
\begin{equation}
  r = \frac{a(1 - e^2)}{1 + e \cos\nu},
\end{equation}
где $a$~--- большая полуось, $e$~--- эксцентриситет, $\nu$~--- истинная
аномалия (угол от направления на перигелий). При $e = 0$ эллипс вырождается
в окружность, а при $e \to 1$ он вытягивается в длинную сигару.

\subsection*{Второй закон (закон площадей)}
Радиус-вектор, проведённый от Солнца к планете, за равные промежутки времени
заметает равные площади:
\begin{equation}
  \frac{dA}{dt} = \frac{1}{2} r^2 \dot{\nu} = \text{const}.
\end{equation}
Это означает, что планета движется быстрее вблизи перигелия и медленнее
вблизи афелия. Как мы увидим позже, второй закон Кеплера есть следствие
сохранения момента импульса.

\subsection*{Третий закон (гармонический закон)}
Квадраты периодов обращения планет относятся как кубы больших полуосей их
орбит:
\begin{equation}
  T^2 = \frac{4\pi^2}{GM} a^3,
\end{equation}
где $M$~--- масса Солнца, $G$~--- гравитационная постоянная. Этот закон
связывает геометрию орбиты (большую полуось $a$) с временной
характеристикой~--- периодом $T$.

\subsection*{Уравнение Кеплера}

Для практического вычисления положения планеты на орбите в заданный момент
времени используется \textbf{уравнение Кеплера}:
\begin{equation}
  M = E - e \sin E,
\end{equation}
где $M = n(t - t_0)$~--- средняя аномалия ($n = 2\pi / T$~--- среднее
движение), $E$~--- эксцентрическая аномалия. Это трансцендентное уравнение
не имеет решения в замкнутой форме и решается численно, например методом
Ньютона. Связь эксцентрической аномалии с истинной аномалией даётся
формулой:
\begin{equation}
  \tan\frac{\nu}{2} = \sqrt{\frac{1+e}{1-e}} \, \tan\frac{E}{2}.
\end{equation}

Законы Кеплера были чисто эмпирическими~--- Кеплер не знал, \emph{почему}
планеты движутся именно так. Объяснение дал Ньютон.

%=============================================================================
\section{Ньютон}
%=============================================================================

Исаак Ньютон (1643--1727) совершил великий синтез, объединив земную механику
Галилея с небесными законами Кеплера в рамках единой теории. Для этого ему
пришлось изобрести два мощнейших математических инструмента~---
дифференциальное и интегральное исчисление (параллельно с Лейбницем).

Ньютон сформулировал \textbf{закон всемирного тяготения}: любые два тела
притягиваются с силой, пропорциональной произведению их масс и обратно
пропорциональной квадрату расстояния между ними:
\begin{equation}
  F = \frac{G M m}{r^2}.
\end{equation}

Из этого единственного закона, вместе со вторым законом Ньютона
$\mathbf{F} = m \mathbf{a}$, можно вывести все три закона Кеплера как
математические следствия. В частности, второй закон Кеплера (закон
площадей) вытекает из того, что гравитационная сила центральна~--- направлена
вдоль радиус-вектора, а значит, момент импульса сохраняется.

Особенно выразительно \textbf{мысленный эксперимент Ньютона с пушкой}.
Представим себе пушку на вершине очень высокой горы. Если выстрелить
горизонтально с малой скоростью, ядро упадёт по параболе~--- это
баллистическое движение Галилея. Если увеличивать скорость, траектория
будет <<загибаться>> за горизонт. При достаточно большой скорости ядро
вообще не упадёт на Землю, а будет двигаться по замкнутому эллипсу~---
это орбитальное движение. При ещё большей скорости (первой космической)
орбита станет круговой. Таким образом, падение яблока и обращение Луны~---
это один и тот же процесс, управляемый одной и той же силой тяготения.

Уравнение движения в гравитационном поле записывается в виде:
\begin{equation}
  m \ddot{\mathbf{r}} = -\frac{G M m}{r^3} \mathbf{r},
\end{equation}
или, сокращая на $m$ и вводя гравитационный параметр $\mu = GM$:
\begin{equation}
  \ddot{\mathbf{r}} = -\frac{\mu}{r^3} \mathbf{r}.
\end{equation}

Это дифференциальное уравнение второго порядка, решениями которого являются
конические сечения: эллипсы (для связанных орбит), параболы (для
параболических траекторий) и гиперболы (для пролётных траекторий). Так
задача двух тел решается аналитически~--- Ньютон показал, что гравитация
порождает в точности те орбиты, которые наблюдал Кеплер.

Энергия системы определяет тип орбиты. Полная удельная энергия:
\begin{equation}
  \mathcal{E} = \frac{v^2}{2} - \frac{\mu}{r} = -\frac{\mu}{2a}.
\end{equation}
При $\mathcal{E} < 0$ орбита эллиптическая, при $\mathcal{E} = 0$~---
параболическая, при $\mathcal{E} > 0$~--- гиперболическая.

%=============================================================================
\section{Задача трёх тел}
%=============================================================================

Если задача двух тел решается аналитически, то уже задача трёх тел не
имеет общего решения в замкнутой форме. Это было одной из величайших
загадок математической физики XVIII--XIX веков. Многие выдающиеся математики
пытались решить эту задачу, и каждый раз обнаруживали непреодолимые
трудности.

Уравнения движения для трёх тел, взаимодействующих гравитационно,
записываются в виде системы:
\begin{equation}
  m_i \ddot{\mathbf{r}}_i = \sum_{\substack{j=1 \\ j \neq i}}^{3}
  \frac{G m_i m_j}{|\mathbf{r}_j - \mathbf{r}_i|^3}
  (\mathbf{r}_j - \mathbf{r}_i), \quad i = 1, 2, 3.
\end{equation}

Это система из 18 обыкновенных дифференциальных уравнений первого порядка
(по три компоненты координат и три компоненты скорости для каждого из трёх
тел). Несмотря на кажущуюся простоту, общее аналитическое решение не
существует.

Фундаментальный вклад в понимание задачи трёх тел внёс Анри Пуанкаре в
конце XIX века. Исследуя ограниченную задачу трёх тел (когда масса одного
из тел пренебрежимо мала), Пуанкаре обнаружил, что траектории обладают
свойством \textbf{чувствительности к начальным условиям}: сколь угодно малое
изменение начальных данных может привести к радикально иному поведению
системы через достаточно большое время. Это было первое открытие того, что
мы сегодня называем \textbf{детерминированным хаосом}.

Пуанкаре показал, что в фазовом пространстве задачи трёх тел существуют
области, в которых траектории переплетаются бесконечно сложным образом.
Это означает, что долгосрочное предсказание движения трёх взаимодействующих
тел в общем случае невозможно, даже если мы знаем точные уравнения движения.

Открытие Пуанкаре имело далеко идущие последствия. Оно показало, что
классический детерминизм Лапласа~--- представление о том, что, зная
начальные условия, можно предсказать будущее на сколь угодно далёкий
срок~--- имеет фундаментальные ограничения даже в рамках классической
механики, безо всякой квантовой неопределённости.

Для практических приложений особое значение имеет \textbf{ограниченная задача
трёх тел}, в которой одно из тел (например, космический аппарат) имеет
пренебрежимо малую массу и не влияет на движение двух других (например,
Солнца и Земли). Именно в этой задаче Лагранж нашёл замечательные точки
равновесия.

%=============================================================================
\section{Точки Лагранжа}
%=============================================================================

Жозеф Луи Лагранж в 1772 году, исследуя ограниченную задачу трёх тел,
обнаружил пять особых точек, в которых малое тело может находиться в
равновесии относительно двух массивных тел. Эти точки получили название
\textbf{точек Лагранжа} (или точек либрации) $L_1, L_2, L_3, L_4, L_5$.

Три точки~--- $L_1$, $L_2$, $L_3$~--- расположены на прямой, соединяющей
два массивных тела (\textbf{коллинеарные} точки). Точка $L_1$ находится между
телами, $L_2$~--- за менее массивным телом, $L_3$~--- за более массивным
телом. Все три коллинеарные точки \textbf{неустойчивы}: малое отклонение
от равновесия нарастает со временем. Однако вокруг них существуют
периодические орбиты (орбиты Лиссажу и гало-орбиты), на которых можно
удерживать космический аппарат с минимальными затратами топлива.

Две оставшиеся точки~--- $L_4$ и $L_5$~--- образуют равносторонние
треугольники с двумя массивными телами (\textbf{треугольные} точки).
Замечательный результат Лагранжа: эти точки \textbf{устойчивы}, если
отношение масс двух массивных тел удовлетворяет условию:
\begin{equation}
  \frac{m_2}{m_1 + m_2} < \frac{1}{2}\left(1 - \sqrt{\frac{23}{27}}\right)
  \approx 0{,}0385,
\end{equation}
что соответствует отношению масс приблизительно $1/25{,}96$. Для системы
Солнце--Юпитер это условие выполняется, и в окрестности точек $L_4$ и $L_5$
Юпитера находятся большие скопления астероидов. Те, что движутся впереди
Юпитера по орбите (вблизи $L_4$), называются \textbf{<<греками>>}
(<<Greeks>>), а те, что позади (вблизи $L_5$)~--- \textbf{<<троянцами>>}
(<<Trojans>>).

Точки Лагранжа широко используются в современной космонавтике:
\begin{itemize}
  \item Точка $L_1$ системы Солнце--Земля~--- идеальное место для
    наблюдения за Солнцем. Здесь расположена солнечная обсерватория
    \textbf{SOHO} (Solar and Heliospheric Observatory).
  \item Точка $L_2$ системы Солнце--Земля~--- превосходная позиция для
    космических телескопов, так как Земля и Солнце всегда находятся с одной
    стороны и не мешают наблюдениям глубокого космоса. Здесь работает
    телескоп \textbf{JWST} (James Webb Space Telescope).
  \item Точка $L_2$ также использовалась для обсерватории Планка и
    космического аппарата Gaia.
\end{itemize}

%=============================================================================
\section{Общая теория относительности}
%=============================================================================

К концу XIX века ньютоновская небесная механика достигла впечатляющих
успехов. Однако оставалась одна упрямая аномалия: прецессия перигелия
Меркурия.

Наблюдения показывали, что перигелий орбиты Меркурия смещается
(прецессирует) со скоростью около $5600''$ (угловых секунд) за столетие.
Бо\'{л}ьшую часть этой прецессии~--- примерно $5557''$/столетие~--- можно
объяснить влиянием других планет (прежде всего Венеры и Юпитера). Однако
оставался необъяснённый остаток: примерно $43''$ за столетие. Это
расхождение было малым, но несомненным, и не поддавалось объяснению в
рамках ньютоновской теории.

Альберт Эйнштейн в 1915 году создал \textbf{общую теорию относительности}
(ОТО), в которой гравитация описывается не как сила, а как искривление
пространства-времени массой. Уравнения Эйнштейна связывают геометрию
пространства-времени (тензор кривизны) с распределением
материи-энергии (тензор энергии-импульса):
\begin{equation}
  R_{\mu\nu} - \frac{1}{2}R\,g_{\mu\nu} = \frac{8\pi G}{c^4}\,T_{\mu\nu}.
\end{equation}

В слабом гравитационном поле ОТО даёт малые поправки к ньютоновской теории.
Для орбитального движения эффективный потенциал приобретает дополнительный
член, и уравнение орбиты (для переменной $u = 1/r$) принимает вид:
\begin{equation}
  \frac{d^2 u}{d\nu^2} + u = \frac{GM}{L^2} + \frac{3GM}{c^2}\,u^2,
\end{equation}
где $L$~--- удельный момент импульса, а последний член $\frac{3GM}{c^2}\,u^2$
представляет собой \textbf{релятивистскую поправку}, отсутствующую в
ньютоновской теории. Именно этот член приводит к дополнительной прецессии
перигелия.

Скорость релятивистской прецессии перигелия за один оборот составляет:
\begin{equation}
  \Delta\varphi = \frac{6\pi G M}{c^2 a (1 - e^2)},
\end{equation}
где $a$~--- большая полуось, $e$~--- эксцентриситет орбиты. Для Меркурия
это даёт как раз $43''$ за столетие~--- в точном согласии с наблюдениями.

Объяснение аномальной прецессии Меркурия стало одним из первых и самых
убедительных подтверждений общей теории относительности. Эйнштейн
впоследствии писал, что когда он увидел, что его уравнения дают правильное
значение прецессии, он испытал сильнейшее волнение.

Сегодня эффекты ОТО необходимо учитывать в работе спутниковых навигационных
систем (GPS, ГЛОНАСС), так как различие хода часов на орбите и на
поверхности Земли (гравитационное замедление времени) составляет десятки
микросекунд в сутки, что приводит к ошибкам позиционирования в километры,
если не вносить соответствующие поправки.

%=============================================================================
\section{Заключение}
%=============================================================================

Мы проследили путь небесной механики длиной более двух тысяч лет~--- от
геоцентрических эпициклов Птолемея до уравнений Эйнштейна и хаоса Пуанкаре.
На этом пути были совершены величайшие научные революции: смена картины мира
(Коперник), создание экспериментальной науки (Галилей), открытие
эмпирических законов (Кеплер), великий синтез и изобретение математического
анализа (Ньютон), обнаружение хаоса в детерминированных системах (Пуанкаре),
пересмотр представлений о пространстве и времени (Эйнштейн).

Но есть существенное отличие между прошлым и настоящим. На протяжении
столетий учёные \emph{наблюдали} за движением естественных небесных тел и
стремились понять законы, которым это движение подчиняется. Это была
классическая небесная механика~--- наука описания и предсказания.

Сегодня акцент сместился: мы не просто наблюдаем за небесными телами, а
\emph{создаём} и \emph{контролируем} искусственные спутники и космические
аппараты. Мы проектируем орбиты, выполняем манёвры, поддерживаем
группировки из сотен и тысяч спутников. Это порождает совершенно новый
пласт математических задач, выходящих за рамки классической небесной
механики:
\begin{itemize}
  \item \textbf{Теория управления}~--- как оптимально перевести спутник
    с одной орбиты на другую, минимизируя расход топлива.
  \item \textbf{Оптимальное управление}~--- задачи Понтрягина, принцип
    максимума, манёвры с малой тягой.
  \item \textbf{Управление группировками}~--- координация сотен спутников
    для обеспечения связи, навигации, дистанционного зондирования Земли.
  \item \textbf{Уклонение от столкновений}~--- с ростом числа объектов на
    орбите проблема космического мусора становится всё более острой.
\end{itemize}

Таким образом, небесная механика не только не утратила своей актуальности,
но, напротив, переживает новый расцвет, обогащаясь методами теории
управления, оптимизации и информатики. История, которую мы сегодня
рассмотрели~--- от Птолемея до спутниковых группировок~--- продолжается, и
самые интересные её главы, возможно, ещё впереди.

\end{document}
